\documentclass[12pt]{article}
\input{preamble}

\pagestyle{fancy}
\fancyhf{}

\rhead{Nørre Gymnasium\\
1.v
}
\cfoot{Side \thepage \hspace{1pt} af \pageref{LastPage}}

%Husk at rette modul og dato!
\lhead{Matematik B\\
Modul 2
}
\chead{08/11/2021
}


\begin{document}

%Udfyld afsnit herunder og lav til egen Latex-fil


%Overskrift
\begin{center}
\Huge
Regneregler
\end{center}
Vi starter med at huske os selv på regnearternes hierarki også kaldet operatorpræcedens, der fortæller os i hvilken rækkefølge, vi skal anvende operatorerne i et givent regnestykke. Altså i hvilken rækkefølge, vi skal lægge sammen, gange, dividere, tage potenser osv. 
\begin{defn}[Regnearternes hierarki]
I en udregning anvender vi operatorerne i følgende rækkefølge:
\begin{enumerate}[label=\roman*)]
\item Parentes. Betegnes med $()$. (Alt, der står i parentes udregnes først efter rækkefølgen bestemt for de resterende operatorer).
\item Fakultet. Betegnes med $!$. Vi husker på, at for $n\in \mathbb{N}$ er $n!$ defineret som 
\begin{align*}
n! = \begin{cases}
1 \ &\textnormal{ for }n=0,\\
n(n-1)(n-2)\cdots 2\cdot 1 \ &\textnormal{ for }n>0.
\end{cases}
\end{align*}
\item Potenser og rødder. Et tal $a$ i $n$'te potens og $n$'te rod betegnes med henholdsvist $a^n$ og $\sqrt[n]{a}$.
\item Multiplikation og division. Betegnes med henholdsvist $\cdot$ og $/$.
\item Addition og subtraktion. Betegnes med henholdsvist $+$ og $-$.
\end{enumerate}
\end{defn} 
\begin{exa}
Lad os betragte regnestykket
\begin{align}\label{eq:exa1}
7+10-\underbrace{(5-2\cdot \frac{3}{6}+3!^2)}_{\textnormal{Parentes 1}}+\underbrace{(7-9)}_{\textnormal{Parentes 2}}*4.
\end{align}
Vi starter med at udregne Parentes 1. Vi følger regnearternes hierarki:
\begin{align*}
(5-2\cdot \frac{3}{6}+3!^2) &\overset{\textnormal{ii)}}{=} (5-2\cdot \frac{3}{6}+6^2)\\
							&\overset{\textnormal{iii)}}{=} (5-2\cdot \frac{3}{6}+36)\\
							&\overset{\textnormal{iv)}}{=} (5-1+36)\\
							&\overset{\textnormal{v)}}{=} (40).
\end{align*}
Og Parentes 2 tilsvarende:
\begin{align*}
(7-9) &\overset{\textnormal{v)}}{=} (-2).
\end{align*}
Disse indsættes nu i \eqref{eq:exa1}, og vi anvender igen regnearternes hierarki til at udregne:
\begin{align*}
7+10+\underbrace{40}_{\textnormal{Par. 1}} + \underbrace{(-2)}_{\textnormal{Par. 2}}\cdot 3 &\overset{\textnormal{iv)}}{=} 7+10+40 -6 \\
&\overset{\textnormal{v)}}{=} 51.
\end{align*}
\end{exa}
\section*{Brøker}
\stepcounter{section}
\begin{enumerate}[label=\roman*)]
\item Addition og subtraktion af brøker. For to brøker $\frac{a}{b}$ og $\frac{c}{d}$ er summen og differensen givet ved
\begin{align*}
\frac{a}{b} \pm \frac{c}{d} = \frac{ad\pm cb}{bd}.
\end{align*}
\item Multiplikation af brøker. Produktet mellem $\frac{a}{b}$ og $\frac{c}{d}$ er givet ved
\begin{align*}
\frac{a}{b}\cdot \frac{c}{d} = \frac{ac}{bd}.
\end{align*}
Specielt er produktet mellem en brøk $\frac{a}{b}$ og et tal $c$ givet ved
\begin{align*}
c\cdot \frac{a}{b} = \frac{ca}{b}.
\end{align*}
\item Division af brøker. Forholdet mellem to brøker $\frac{a}{b}$ og $\frac{c}{d}$ er givet ved
\begin{align*}
\frac{\frac{a}{b}}{\frac{c}{d}} = \frac{ad}{bc}.
\end{align*}
(Vi ganger med den omvendte brøk.) Specielt er en brøk $\frac{a}{b}$ divideret med et tal $c$ givet ved
\begin{align*}
\frac{\frac{a}{b}}{c} = \frac{a}{bc},
\end{align*}
og et tal $c$ divideret med en brøk $\frac{a}{b}$ givet ved
\begin{align*}
\frac{c}{\frac{a}{b}} = \frac{ac}{b}
\end{align*}
\item Brøker og potenser/rødder. En brøk $\frac{a}{b}$ opløftet i et tal $c$ er givet ved 
\begin{align*}
\left(\frac{a}{b}\right)^{c} = \frac{a^c}{b^c}.
\end{align*}
$n$'te roden af en brøk $\frac{a}{b}$ er givet ved 
\begin{align*}
\sqrt[n]{\frac{a}{b}} = \frac{\sqrt[n]{a}}{\sqrt[n]{b}}.
\end{align*}
\end{enumerate}
\section*{Potenser/rødder og det udvidede potensbegreb}
\stepcounter{section}

Når vi opløfter et tal $a\in \mathbb{R}$ i et naturligt tal $n\in \mathbb{N}$ så tilsvarer det
\begin{align}\label{eq:powerdef}
a^n = \underbrace{a\cdot a \cdots a}_{n\textnormal{ gange}},
\end{align}
vi multiplicerer altså tallet $a$ $n$ gange med sig selv. Vi vil snart forklare, hvad det betyder at opløfte et tal i både en positiv og negativ brøk. Vi vil først diskutere regnereglerne for potenser. Vi vil udnytte repræsentationen \eqref{eq:powerdef} for at udlede regnereglerne for multiplikation

\end{document}